%!TEX root = ../dissertation.tex
\chapter{Related Work}\label{chapt:related}

% =====================================================================
%  DRAFT v1 (Asma) -- full prose. ~10 pages target.
%  HIGHLIGHT markers = citations to verify / optional additions.
% =====================================================================

This chapter situates the thesis within the literature.
Section~\ref{sect:rel_supervised} reviews supervised composed image retrieval
and its benchmarks. Section~\ref{sect:rel_zeroshot} surveys zero-shot CIR.
Section~\ref{sect:rel_icir} introduces instance-level retrieval and the i-CIR
benchmark. Section~\ref{sect:rel_caption} discusses caption- and LLM-based
approaches that, like ours, route the query through language.
Section~\ref{sect:rel_backbones} reviews the vision--language backbones used
for retrieval. The chapter closes with the gap this thesis addresses
(Section~\ref{sect:rel_gap}).

% =====================================================================
\section{Supervised Composed Image Retrieval}\label{sect:rel_supervised}
% =====================================================================

The composed image retrieval task was popularised by TIRG (Text Image Residual
Gating)~\cite{vo2019tirg}, which framed retrieval with a (reference image,
modification text) query and proposed an \emph{early-fusion} network that
combines the image and text features into a single composed embedding via a
gated residual connection. TIRG established the now-standard training recipe:
supervise a fusion module with a metric-learning loss on annotated triplets so
that the composed query embedding is pulled towards the target-image embedding.

Subsequent work scaled the task to more realistic data and introduced the
benchmarks that the field still uses. \emph{Fashion-IQ}~\cite{wu2021fashioniq}
targets fashion product retrieval with natural-language relative feedback,
while \emph{CIRR}~\cite{liu2021cirr} moved to open-domain real-life images and
explicitly highlighted the need for models that reason about fine-grained
visual differences rather than matching on global appearance. A recurring
limitation of these benchmarks is that relevance is judged at the
\emph{category} or single-target level, and the modification language is often
correlated with easy visual cues, which inflates measured performance -- an
observation that directly motivates the instance-level protocol discussed in
Section~\ref{sect:rel_icir}.

% =====================================================================
\section{Zero-Shot Composed Image Retrieval}\label{sect:rel_zeroshot}
% =====================================================================

Because triplet supervision is expensive and does not transfer well across
domains, a line of work removes CIR-specific training entirely and builds on
frozen vision--language models. \emph{Pic2Word}~\cite{saito2023pic2word}
introduced \emph{textual inversion} for CIR: it learns, using only image data
(no triplets), a mapping that converts a reference image into a single
pseudo-word token; the modification text is then written as a prompt containing
this token (e.g. ``a photo of $[\ast]$, but red''), and the composed prompt is
embedded by the frozen CLIP text encoder. SEARLE and the CIRCO
benchmark~\cite{baldrati2023circo} refined this idea, optimising the inverted
token to be close to a set of descriptive words and providing a CIR test set
with multiple validated ground-truth targets to reduce annotation noise.
LinCIR~\cite{gu2024lincir} pushed efficiency further by training the inversion
using \emph{language only}, without any images.

These textual-inversion methods share the philosophy of this thesis -- express
the composed query as \emph{text} so that the frozen text encoder can handle it
-- but differ in mechanism. They compress the reference image into a single
opaque token and still rely on the user's short modification phrase as the rest
of the prompt. In contrast, this thesis uses a generative MLLM to produce a
\emph{full, human-readable description of the imagined target}, which captures
the composition explicitly rather than relegating the reference to one learned
token.

% =====================================================================
\section{Instance-Level Retrieval and the i-CIR Benchmark}\label{sect:rel_icir}
% =====================================================================

Instance-level retrieval -- finding images of a \emph{specific} object instance
rather than its category -- has a long tradition in classical image
retrieval~\cite{philbin2007oxford}, where evaluation pools deliberately contain
visually similar distractors of the same category. The i-CIR
benchmark~\cite{psomas2025icir} brings this rigour to composed image retrieval.
In i-CIR a retrieval is correct only if it returns the \emph{same instance}
modified as the query specifies, and each query's evaluation is conducted
against a pool engineered to contain hard negatives: the same instance under a
different modification, visually similar but text-violating images, and
text-satisfying images of a different instance. This design penalises models
that exploit only one modality and demands genuine composition.

Alongside the benchmark, i-CIR contributes a strong zero-shot baseline,
\textsc{Basic} (a ``Baseline Approach for SurprIsingly strong Composition'').
\textsc{Basic} combines a frozen vision--language backbone with a sequence of
training-free post-processing steps -- feature centering, dimensionality
reduction, query expansion, and a \emph{contextualization} of the modification
text that prepends or appends a subject term to the raw phrase -- and fuses the
image and text similarities into a final score. \textsc{Basic} demonstrates
that careful use of frozen features is a deceptively strong baseline, and it is
the method this thesis builds upon and compares against. The present work keeps
the late-fusion philosophy of \textsc{Basic} but argues that an MLLM-generated
target caption is a more effective way to inject the textual signal than
contextualizing the bare modification phrase.

% =====================================================================
\section{Caption- and LLM-Based Composed Retrieval}\label{sect:rel_caption}
% =====================================================================

A complementary family of zero-shot methods routes the query through
\emph{generated language}. CIReVL (Compositional Image Retrieval via
Vision-by-Language)~\cite{karthik2024cirevl} is the closest in spirit to this
thesis: it is training-free and uses a captioner to describe the reference
image, then prompts a large language model to rewrite that caption according to
the modification, yielding a target description that is embedded by the frozen
CLIP text encoder. This confirms the central intuition -- that an explicit
textual description of the intended target is a powerful zero-shot query -- but
CIReVL operates in two stages (caption the reference, then edit the caption with
a separate LLM) and was evaluated on category-level benchmarks.

This thesis differs in three respects. First, it uses a single
\emph{multimodal} model that sees the reference image and the modification
\emph{together} and writes the target caption in one step, rather than
captioning and editing in separate text-only stages. Second, it does not
\emph{replace} the image and text signals with the caption; it \emph{adds} the
caption as a third modality and fuses all three by a clamped product, so the
visual identity of the instance is preserved -- which is essential at the
instance level. Third, it is evaluated on the demanding i-CIR instance-level
benchmark and on the full 752K-image database, rather than on category-level
test sets.

% =====================================================================
\section{Vision--Language Backbones for Retrieval}\label{sect:rel_backbones}
% =====================================================================

The choice of frozen backbone strongly determines zero-shot retrieval quality.
CLIP~\cite{radford2021clip} and its open reproductions trained on
LAION~\cite{ilharco2021openclip,schuhmann2022laion5b} remain the default
encoders for zero-shot CIR. SigLIP~\cite{zhai2023siglip} and
SigLIP2~\cite{tschannen2025siglip2} provide more strongly aligned image--text
spaces through the sigmoid objective and improved training, and are
increasingly adopted for retrieval. This thesis compares CLIP ViT-L/14, CLIP
ViT-H/14, SigLIP ViT-L/16, and SigLIP2 g/16-384 on the full i-CIR benchmark
under an identical pipeline, isolating the effect of the backbone
(Section~\ref{sect:res_backbones}).

% =====================================================================
\section{Positioning and Gap}\label{sect:rel_gap}
% =====================================================================

In summary, supervised CIR requires expensive triplet annotation and transfers
poorly; zero-shot textual-inversion methods express the query as text but
compress the reference into a single learned token and still depend on the bare
modification phrase; and the strongest zero-shot baseline on the new
instance-level benchmark, \textsc{Basic}, injects the textual signal by
contextualizing the modification rather than describing the target. The gap
this thesis addresses is therefore the following: \emph{in the zero-shot,
instance-level setting, can an MLLM-generated description of the imagined
target, fused as a third modality with the image and text signals, provide a
stronger and simpler textual signal than the elaborate preprocessing of
\textsc{Basic}?} The remainder of the thesis answers this question
affirmatively.
